(1) With $X_p(x) = (1, x, \tfrac12 Q_p(x))$ the incidence condition reads $Q_p(x) = 2 \cdot 1 \cdot \tfrac12 Q_p(x)$, an identity. Applying $\iota$ gives $(Q_p(x), x, \tfrac12)$, whose incidence check is $2 \cdot Q_p(x) \cdot \tfrac12 = Q_p(x)$, again an identity; and $\iota^2(s,x,t) = \iota(2t, x, \tfrac12 s) = (s, x, t)$.

(2) $\iota(X_p(x)) = (Q_p(x), x, \tfrac12)$, so $\dsc$ returns $x/Q_p(x)$, which is the classical reciprocal $\rho(x)$ of \cite[Lem.~4.23(2)]{ButlerT1}.

(3) For the stated $x$, $Q_p(x) = f \cdot 1 - f^{-1} \cdot f^2 = f - f = 0$, so $X_p(x) = (1, x, 0)$ and $\iota(X_p(x)) = (0, x, \tfrac12)$: the first slot vanishes, which is the defining condition of an ideal point, and $\dsc$ requires division by it. The value is $[\,x + \tfrac12 \nInf\,]$, exactly the ideal value of \cite[Thm.~4.10(3)]{ButlerT1}.

(4) With $t_{\mathrm{PG}} = t_s + \int \beta f^{-1}dr$ and $\beta^{2} = 2M/r$, $dt_{\mathrm{PG}} = dt_s + (\beta/f)dr$, so $k = \sqrt{2M}\,dt_s + \sqrt r\,(1 + 2M/rf)\,dr = \sqrt{2M}\,dt_s + (\sqrt r/f)\,dr$, using $f + 2M/r = 1$. In the static cometric $(g_p^{-1})^{tt} = 1/f$, $(g_p^{-1})^{rr} = -f$, so $Q^*_p(\xi) = 2M/f - f\,(r/f^{2}) = (2M - r)/f = -r$, an identity of rational functions in $r$ with no pole at the horizon; in the Painlev\'e-Gullstrand frame of \cite[Thm.~8.7]{ButlerT4} the same value is read off the polynomial Gram matrix as $g_k = (m^{2} - r) - 2m^{2} + m^{2} = -r$ with $m = \sqrt{2M}$, which is that theorem's clause (3). Hence $\iota(X_p(\xi)) = (-r, \xi, \tfrac12)$ has nonvanishing first slot at every $r > 0$, where $\dsc$ returns $-\xi/r$, and first slot tending to $0$ as $r \to 0^+$, which is the ideal stratum. (The Killing-time covector $dt_s$ itself, with $Q^*_p(dt_s) = 1/f$, spans the radical line only on $E$ \cite[Thm.~8.7(4)]{ButlerT4}, and its static norm has a pole at the horizon where nothing degenerates; $k$ is the regular generator.) One further line closes the junction with Step 1, computed in the Painlev\'e-Gullstrand chart where Propositions \claimref{singularity.cones-fold} and \claimref{singularity.degeneration-exactly} fold the cones: the PG cometric has $g^{tt} = 1$, $g^{tr} = -\beta$, $g^{rr} = \beta^{2} - 1$, so for $k = \sqrt{2M}\,dt_{\mathrm{PG}} + \sqrt r\,dr$ the dual's $t$-component is $\sqrt{2M} - \beta\sqrt r = 0$ identically and its $r$-component is $-\beta\sqrt{2M} + (\beta^{2} - 1)\sqrt r = -\sqrt r$: $k^{\sharp} = -\sqrt r\,\er$ exactly, at every radius and with no horizon pole, in the same frame as Step 1, so the corank-one radical of the completion and Step 1's dying direction $\er$ are the same direction outright, not merely in the limit. That the failure locus is exactly $E$, and that no shadow wall forms elsewhere, is \cite[Thm.~8.12]{ButlerT4}. All identities in (1) to (4) are machine-verified in exact symbolic arithmetic (Section \ref{sec:verification}, parts L1 to L14), and the nullity of the lift, $X_p(x)^{2} = 0$ in $\Cl(2,4)$ with the stated normalisation, blade-exactly.
